\documentclass[11pt]{beamer}
\usepackage[utf8]{inputenc}
\usepackage[spanish]{babel}
\usepackage{amsmath}
\usepackage{amssymb}

\usetheme{Madrid}
\usecolortheme{seahorse}

\title{Problema 4: Perfiles Clínicos}
\subtitle{Operaciones de Conjuntos}
\author{}
\date{}

\begin{document}
	
	\begin{frame}
		\titlepage
	\end{frame}
	
	%--- FRAME 1: Enunciado ---
	\begin{frame}{Descripción del Problema}
		Se estudia un grupo de \textbf{20 personas} para identificar perfiles clínicos.
		\[
		U = \{1, 2, 3, \ldots, 20\}
		\]
		\begin{itemize}
			\item \textbf{Ansiedad elevada} $(A)$: altos niveles de activación emocional.
			\item \textbf{Dificultades sociales} $(S)$: problemas en habilidades e interacción.
			\item \textbf{Estrés crónico} $(E)$: sobrecarga sostenida o fatiga emocional.
		\end{itemize}
	\end{frame}
	
	%--- FRAME 2: Conjuntos ---
	\begin{frame}{Conjuntos Dados}
		\begin{block}{Conjuntos por extensión}
			\begin{align*}
				A &= \{2, 4, 6, 8, 10, 12, 14, 16, 18, 20\} \\[6pt]
				S &= \{5, 7, 9, 11, 13\} \\[6pt]
				E &= \{3, 6, 9, 12, 15, 18\}
			\end{align*}
		\end{block}
		\begin{alertblock}{Observación clave}
			$A \cap S = \emptyset$ \quad (ningún elemento de $S$ pertenece a $A$)
		\end{alertblock}
	\end{frame}
	
	%--- FRAME 3: Parte a) - enunciado ---
	\begin{frame}{Parte a) — Conjunto por Comprensión}
		\begin{block}{Se pide determinar por comprensión:}
			\[
			(A - E) \cup (S \cap E)
			\]
		\end{block}
		\bigskip
		Resolvemos cada operación por separado:
		\begin{enumerate}
			\item Calcular $A - E$
			\item Calcular $S \cap E$
			\item Unir ambos resultados
		\end{enumerate}
	\end{frame}
	
	%--- FRAME 4: Parte a) - desarrollo ---
	\begin{frame}{Parte a) — Desarrollo}
		\textbf{Paso 1:} $A - E$ (elementos de $A$ que \textbf{no} están en $E$)
		\[
		A - E = \{2, 4, \cancel{6}, 8, 10, \cancel{12}, 14, 16, \cancel{18}, 20\} = \{2, 4, 8, 10, 14, 16, 20\}
		\]
		
		\medskip
		\textbf{Paso 2:} $S \cap E$ (elementos en $S$ \textbf{y} en $E$)
		\[
		S \cap E = \{5, 7, \mathbf{9}, 11, 13\} \cap \{3, 6, \mathbf{9}, 12, 15, 18\} = \{9\}
		\]
		
		\medskip
		\textbf{Paso 3:} Unión
		\[
		(A - E) \cup (S \cap E) = \{2, 4, 8, 9, 10, 14, 16, 20\}
		\]
	\end{frame}
	
	%--- FRAME 5: Parte a) - resultado y comprensión ---
	\begin{frame}{Parte a) — Resultado}
		\begin{exampleblock}{Resultado}
			\[
			(A - E) \cup (S \cap E) = \{2,\, 4,\, 8,\, 9,\, 10,\, 14,\, 16,\, 20\}
			\]
		\end{exampleblock}
		
		\bigskip
		\textbf{Descripción por comprensión:}
		
		\medskip
		Conjunto de personas que \textbf{presentan ansiedad elevada pero no estrés crónico},
		o que \textbf{tienen simultáneamente dificultades sociales y estrés crónico}.
	\end{frame}
	
	%--- FRAME 6: Parte b) - enunciado ---
	\begin{frame}{Parte b) — Conjunto por Extensión}
		\begin{block}{Se pide escribir por extensión y determinar la cardinalidad de:}
			\[
			(A - S)^c \cap A^c - E
			\]
		\end{block}
		
		\bigskip
		Aplicando precedencia estándar de operaciones:
		\[
		\bigl[(A - S)^c \cap A^c\bigr] - E
		\]
		
		\medskip
		Resolvemos paso a paso:
		\begin{enumerate}
			\item Calcular $A - S$
			\item Calcular $(A-S)^c$
			\item Intersectar con $A^c$
			\item Restar $E$
		\end{enumerate}
	\end{frame}
	
	%--- FRAME 7: Parte b) - desarrollo 1 ---
	\begin{frame}{Parte b) — Desarrollo (1/2)}
		\textbf{Paso 1:} $A - S$
		
		Como $A \cap S = \emptyset$, entonces:
		\[
		A - S = A = \{2, 4, 6, 8, 10, 12, 14, 16, 18, 20\}
		\]
		
		\medskip
		\textbf{Paso 2:} $(A - S)^c = A^c$
		\[
		A^c = U - A = \{1, 3, 5, 7, 9, 11, 13, 15, 17, 19\}
		\]
		
		\medskip
		\textbf{Paso 3:} $(A-S)^c \cap A^c = A^c \cap A^c$
		\[
		A^c \cap A^c = A^c = \{1, 3, 5, 7, 9, 11, 13, 15, 17, 19\}
		\]
	\end{frame}
	
	%--- FRAME 8: Parte b) - desarrollo 2 ---
	\begin{frame}{Parte b) — Desarrollo (2/2)}
		\textbf{Paso 4:} $A^c - E$ (elementos de $A^c$ que \textbf{no} están en $E$)
		
		\begin{align*}
			A^c &= \{1, \cancel{3}, 5, 7, \cancel{9}, 11, 13, \cancel{15}, 17, 19\} \\
			E   &= \{3, 6, 9, 12, 15, 18\}
		\end{align*}
		
		\[
		A^c - E = \{1, 5, 7, 11, 13, 17, 19\}
		\]
	\end{frame}
	
	%--- FRAME 9: Parte b) - resultado ---
	\begin{frame}{Parte b) — Resultado}
		\begin{exampleblock}{Resultado por extensión}
			\[
			(A-S)^c \cap A^c - E = \{1,\, 5,\, 7,\, 11,\, 13,\, 17,\, 19\}
			\]
		\end{exampleblock}
		
		\begin{block}{Cardinalidad}
			\[
			\bigl|(A-S)^c \cap A^c - E\bigr| = \mathbf{7}
			\]
		\end{block}
	\end{frame}
	
	%--- FRAME 10: Resumen ---
	\begin{frame}{Resumen de Resultados}
		\begin{center}
			\renewcommand{\arraystretch}{1.8}
			\begin{tabular}{|c|l|c|}
				\hline
				\textbf{Parte} & \textbf{Conjunto} & \textbf{Cardinalidad} \\
				\hline
				a) & $\{2, 4, 8, 9, 10, 14, 16, 20\}$ & $8$ \\
				\hline
				b) & $\{1, 5, 7, 11, 13, 17, 19\}$ & $7$ \\
				\hline
			\end{tabular}
		\end{center}
	\end{frame}
	
\end{document}